% Part: first-order-logic
% Chapter: syntax-and-semantics
% Section: satisfaction

\documentclass[../../../include/open-logic-section]{subfiles}

\begin{document}

\olfileid{fol}{syn}{sat}

\olsection{إرضاء \article{formula} \printtoken{S}{formula}
  في \article{structure} \printtoken{S}{structure}}

\begin{explain}
يرتبط باب التعبيرات، من حدود وال!!{formula}s، بباب
ال!!{structure}s من جهتين: \emph{!!{value}} الحد، و\emph{إرضاء}
!!a{formula}. وتقريب الأولى أن !!{value} الحد !!a{element} من
!!a{structure}؛ فإن كان الحد رمزًا ثابتًا فحسب، فـ!!{value} هي الشيء
الذي عيّنته ال!!{structure} لذلك الثابت. وإن تألّف من
ال!!{function}s، حُسبت !!{value} له من !!{value}s الثوابت ومن
الدوال المعيّنة لرموز الدوال الواردة فيه. وأما !!^a{formula} فتكون
\emph{مُرضاة} في !!a{structure} متى جعل تأويل المحمولات
ال!!{formula} صادقة في مجال ال!!{structure}. وهذا أيضًا يُضبط
بالاستقراء: فتعيينات ال!!{structure} تبيّن مباشرة متى تُرضى
ال!!{formula}s الذرية؛ وأما !!{formula} المركبة، فيتعلّق إرضاؤها
بالرابط الرئيس أو المكمم، وبإرضاء ال!!{subformula}s المباشرة أو
عدم إرضائها.

غير أن في المكممات مزيد نظر؛ إذ قد يكون في ال!!{subformula}
المباشرة ل!!{formula} مكممة !!{variable} حر، وال!!{structure}s
لا تعيّن !!{value}s ال!!{variable}s. فلا بد لمعالجة ذلك من
\emph{إسنادات المتغيرات}. وحينئذ لا نعتبر الإرضاء بالنسبة إلى
!!a{structure} مجردة، بل إلى !!a{structure} ومعها إسناد
ل!!a{variable}.
\end{explain}

\begin{defn}[إسناد المتغيرات]
\emph{إسناد المتغيرات}~${\OLId{latin-ordinary}{s}}$ ل!!a{structure}~$\Struct{M}$ دالة ترسل
كل !!{variable} إلى !!a{element} من~$\Domain {\OLId{latin-looped}{M}}$؛ أي
${\OLId{latin-ordinary}{s}}\colon \Var \to \Domain {\OLId{latin-looped}{M}}$.
\end{defn}

\begin{explain}
في !!^a{structure} تتعيّن !!a{value} كل !!{constant}، وبإسناد
المتغيرات تتعيّن قيمة كل متغير. ونريد، فوق ذلك، أن تدل الحدود
المركبة منهما على !!{element}s من ال!!{domain}. فلذلك نعرّف
!!{value} الحدود بالاستقراء: قيمة ال!!{constant}s والمتغيرات
مأخوذة مما عيّنته ال!!{structure} أو الإسناد، وقيمة الحدود
المركبة تُحسب عوديًا بالدوال التي عيّنتها ال!!{structure}
لل!!{function}s.
\end{explain}

\begin{defn}[!!^{value} الحدود]
إذا كان ${\OLId{latin-ordinary}{t}}$ حدًا من اللغة~$\Lang L$، وكانت $\Struct M$
!!{structure} لـ~$\Lang L$، وكان ${\OLId{latin-ordinary}{s}}$ إسنادًا ل!!a{variable}
في~$\Struct M$، عُرّفت \emph{!!{value}}~$\Value{{\OLId{latin-ordinary}{t}}}{{\OLId{latin-looped}{M}}}[{\OLId{latin-ordinary}{s}}]$ كما يأتي:
\begin{enumerate}
\item \indcase{t}{c}{$\Value{\indfrm}{{\OLId{latin-looped}{M}}}[{\OLId{latin-ordinary}{s}}] = \Assign{\indcomplex}{{\OLId{latin-looped}{M}}}$.}
\item \indcase{t}{x}{$\Value{\indfrm}{{\OLId{latin-looped}{M}}}[{\OLId{latin-ordinary}{s}}] = {\OLId{latin-ordinary}{s}}(\indcomplex)$.}
\item \indcase{t}{\Atom{f}{t_1, \ldots, t_n}}{
\[
\Value{\indfrm}{{\OLId{latin-looped}{M}}}[{\OLId{latin-ordinary}{s}}] = \Assign{{\OLId{latin-ordinary}{f}}}{{\OLId{latin-looped}{M}}}(\Value{{\OLId{latin-ordinary}{t}}_{\OLMathNumeral{1}}}{{\OLId{latin-looped}{M}}}[{\OLId{latin-ordinary}{s}}], \ldots,
\Value{{\OLId{latin-ordinary}{t}}_{\OLId{latin-ordinary}{n}}}{{\OLId{latin-looped}{M}}}[{\OLId{latin-ordinary}{s}}]).
\]}
\end{enumerate}
\end{defn}

\begin{defn}[مغاير بالنسبة إلى~${\OLId{latin-ordinary}{x}}$]
إذا كان ${\OLId{latin-ordinary}{s}}$ إسنادًا ل!!a{variable} في !!a{structure}~$\Struct M$،
فكل إسناد ل!!{variable}~${\OLId{latin-ordinary}{s}}'$ في~$\Struct M$ لا يختلف عن~${\OLId{latin-ordinary}{s}}$ إلا،
على الأكثر، فيما يسنده إلى~${\OLId{latin-ordinary}{x}}$ يسمى \emph{مغايرًا بالنسبة إلى~${\OLId{latin-ordinary}{x}}$}
لـ~${\OLId{latin-ordinary}{s}}$. وإذا كان ${\OLId{latin-ordinary}{s}}'$ مغايرًا بالنسبة إلى~${\OLId{latin-ordinary}{x}}$ لـ~${\OLId{latin-ordinary}{s}}$ كتبنا
$\varAssign{{\OLId{latin-ordinary}{s}}'}{{\OLId{latin-ordinary}{s}}}{{\OLId{latin-ordinary}{x}}}$.
\end{defn}

\begin{explain}
وليس المغاير بالنسبة إلى~${\OLId{latin-ordinary}{x}}$ لإسناد~${\OLId{latin-ordinary}{s}}$ \emph{ملزمًا} بأن يجعل
قيمة~${\OLId{latin-ordinary}{x}}$ غير قيمتها الأولى؛ فكل إسناد مغاير بالنسبة
إلى~${\OLId{latin-ordinary}{x}}$ لنفسه أيضًا.
\end{explain}

\begin{defn}
  إذا كان ${\OLId{latin-ordinary}{s}}$ إسنادًا ل!!a{variable} في !!a{structure}~$\Struct M$
  وكان ${\OLId{latin-ordinary}{m}} \in \Domain{{\OLId{latin-looped}{M}}}$، فإن الإسناد~$\Subst{{\OLId{latin-ordinary}{s}}}{{\OLId{latin-ordinary}{m}}}{{\OLId{latin-ordinary}{x}}}$ هو إسناد
  المتغيرات المعرّف بـ
  \[\Subst{{\OLId{latin-ordinary}{s}}}{{\OLId{latin-ordinary}{m}}}{{\OLId{latin-ordinary}{x}}}({\OLId{latin-ordinary}{y}}) = \begin{cases}
    {\OLId{latin-ordinary}{m}} & \text{إذا كان } {\OLId{latin-ordinary}{y}} \ident {\OLId{latin-ordinary}{x}}\\
    {\OLId{latin-ordinary}{s}}({\OLId{latin-ordinary}{y}}) & \text{فيما عدا ذلك}.
  \end{cases}\]
\end{defn}

فـ$\Subst{{\OLId{latin-ordinary}{s}}}{{\OLId{latin-ordinary}{m}}}{{\OLId{latin-ordinary}{x}}}$ مغاير معيّن بالنسبة إلى~${\OLId{latin-ordinary}{x}}$
لـ~${\OLId{latin-ordinary}{s}}$: يضع !!{element} المجال~${\OLId{latin-ordinary}{m}}$ قيمةً لـ~${\OLId{latin-ordinary}{x}}$، ويبقي قيم
ال!!{variable}s غير~${\OLId{latin-ordinary}{x}}$ على ما عيّنه لها~${\OLId{latin-ordinary}{s}}$.

\begin{defn}[الإرضاء]
\ollabel{defn:satisfaction}
يُعرّف إرضاء !!a{formula}~$!A$ في !!a{structure}~$\Struct M$ نسبة
إلى إسناد ل!!a{variable}~${\OLId{latin-ordinary}{s}}$، ويرمز إليه
$\Sat{{\OLId{latin-looped}{M}}}{!A}[{\OLId{latin-ordinary}{s}}]$، عوديًا كما يأتي. (ونكتب
$\Sat/{{\OLId{latin-looped}{M}}}{!A}[{\OLId{latin-ordinary}{s}}]$ بمعنى «ليس $\Sat{{\OLId{latin-looped}{M}}}{!A}[{\OLId{latin-ordinary}{s}}]$».)
\begin{enumerate}
\tagitem{prvFalse}{%
  \indcase{!A}{\lfalse}{$\Sat/{{\OLId{latin-looped}{M}}}{\indfrm}[{\OLId{latin-ordinary}{s}}]$.}}{}

\tagitem{prvTrue}{%
  \indcase{!A}{\ltrue}{$\Sat{{\OLId{latin-looped}{M}}}{\indfrm}[{\OLId{latin-ordinary}{s}}]$.}}{}

\item \indcase{!A}{\Atom{R}{t_1, \dots, t_n}}{$\Sat{{\OLId{latin-looped}{M}}}{\indfrm}[{\OLId{latin-ordinary}{s}}]$
  إذا وفقط إذا كان $\langle \Value{{\OLId{latin-ordinary}{t}}_{\OLMathNumeral{1}}}{{\OLId{latin-looped}{M}}}[{\OLId{latin-ordinary}{s}}], \dots,
  \Value{{\OLId{latin-ordinary}{t}}_{\OLId{latin-ordinary}{n}}}{{\OLId{latin-looped}{M}}}[{\OLId{latin-ordinary}{s}}] \rangle \in \Assign{{\OLId{latin-looped}{R}}}{{\OLId{latin-looped}{M}}}$.}

\item \indcase{!A}{\eq[t_1][t_2]}{$\Sat{{\OLId{latin-looped}{M}}}{\indfrm}[{\OLId{latin-ordinary}{s}}]$ إذا وفقط
  إذا كان $\Value{{\OLId{latin-ordinary}{t}}_{\OLMathNumeral{1}}}{{\OLId{latin-looped}{M}}}[{\OLId{latin-ordinary}{s}}] = \Value{{\OLId{latin-ordinary}{t}}_{\OLMathNumeral{2}}}{{\OLId{latin-looped}{M}}}[{\OLId{latin-ordinary}{s}}]$.}

\tagitem{prvNot}{%
  \indcase{!A}{\lnot !B}{$\Sat{{\OLId{latin-looped}{M}}}{\indfrm}[{\OLId{latin-ordinary}{s}}]$ إذا وفقط إذا كان
    $\Sat/{{\OLId{latin-looped}{M}}}{!B}[{\OLId{latin-ordinary}{s}}]$.}}{}

\tagitem{prvAnd}{%
  \indcase{!A}{(!B \land !C)}{$\Sat{{\OLId{latin-looped}{M}}}{\indfrm}[{\OLId{latin-ordinary}{s}}]$ إذا وفقط إذا
    كان $\Sat{{\OLId{latin-looped}{M}}}{!B}[{\OLId{latin-ordinary}{s}}]$ و$\Sat{{\OLId{latin-looped}{M}}}{!C}[{\OLId{latin-ordinary}{s}}]$.}}{}

\tagitem{prvOr}{%
  \indcase{!A}{(!B \lor !C)}{$\Sat{{\OLId{latin-looped}{M}}}{\indfrm}[{\OLId{latin-ordinary}{s}}]$ إذا وفقط إذا
    كان $\Sat{{\OLId{latin-looped}{M}}}{!B}[{\OLId{latin-ordinary}{s}}]$ أو $\Sat{{\OLId{latin-looped}{M}}}{!C}[{\OLId{latin-ordinary}{s}}]$ (أو كلاهما).}}{}

\tagitem{prvIf}{%
  \indcase{!A}{(!B \lif !C)}{$\Sat{{\OLId{latin-looped}{M}}}{\indfrm}[{\OLId{latin-ordinary}{s}}]$ إذا وفقط إذا
    كان $\Sat/{{\OLId{latin-looped}{M}}}{!B}[{\OLId{latin-ordinary}{s}}]$ أو $\Sat{{\OLId{latin-looped}{M}}}{!C}[{\OLId{latin-ordinary}{s}}]$ (أو كلاهما).}}{}

\tagitem{prvIff}{%
  \indcase{!A}{(!B \liff !C)}{$\Sat{{\OLId{latin-looped}{M}}}{\indfrm}[{\OLId{latin-ordinary}{s}}]$ إذا وفقط إذا
    كان إما كل من $\Sat{{\OLId{latin-looped}{M}}}{!B}[{\OLId{latin-ordinary}{s}}]$ و$\Sat{{\OLId{latin-looped}{M}}}{!C}[{\OLId{latin-ordinary}{s}}]$، وإما لا
    $\Sat{{\OLId{latin-looped}{M}}}{!B}[{\OLId{latin-ordinary}{s}}]$ ولا $\Sat{{\OLId{latin-looped}{M}}}{!C}[{\OLId{latin-ordinary}{s}}]$.}}{}

\tagitem{prvAll}{%
  \indcase{!A}{\lforall[x][!B]}{$\Sat{{\OLId{latin-looped}{M}}}{\indfrm}[{\OLId{latin-ordinary}{s}}]$ إذا وفقط إذا
    كان، لكل !!{element}~${\OLId{latin-ordinary}{m}} \in \Domain {\OLId{latin-looped}{M}}$،
    $\Sat{{\OLId{latin-looped}{M}}}{!B}[\Subst{{\OLId{latin-ordinary}{s}}}{{\OLId{latin-ordinary}{m}}}{{\OLId{latin-ordinary}{x}}}]$.}}{}

\tagitem{prvEx}{%
  \indcase{!A}{\lexists[x][!B]}{$\Sat{{\OLId{latin-looped}{M}}}{\indfrm}[{\OLId{latin-ordinary}{s}}]$ إذا وفقط إذا
  كان، ل!!{element} واحد على الأقل~${\OLId{latin-ordinary}{m}} \in \Domain {\OLId{latin-looped}{M}}$،
  $\Sat{{\OLId{latin-looped}{M}}}{!B}[\Subst{{\OLId{latin-ordinary}{s}}}{{\OLId{latin-ordinary}{m}}}{{\OLId{latin-ordinary}{x}}}]$.}}{}
\end{enumerate}
\end{defn}

\begin{explain}
ويتبيّن موضع الحاجة إلى إسنادات المتغيرات في
\iftag{notprvEx,notprvAll}{البند الأخير}{البندين الأخيرين}.\iftag{prvAll}{
فلا يصح أن نحدّ إرضاء $\lforall[{\OLId{latin-ordinary}{x}}][!B({\OLId{latin-ordinary}{x}})]$ بقولنا: «لكل ${\OLId{latin-ordinary}{m}} \in
\Domain{{\OLId{latin-looped}{M}}}$، $\Sat{{\OLId{latin-looped}{M}}}{!B({\OLId{latin-ordinary}{m}})}$».}{}\iftag{prvEx}{ ولا أن نحدّ إرضاء
$\lexists[{\OLId{latin-ordinary}{x}}][!B({\OLId{latin-ordinary}{x}})]$ بقولنا: «لـ~${\OLId{latin-ordinary}{m}} \in \Domain{{\OLId{latin-looped}{M}}}$ واحد على
الأقل، $\Sat{{\OLId{latin-looped}{M}}}{!B({\OLId{latin-ordinary}{m}})}$».}{} وذلك أن ${\OLId{latin-ordinary}{m}} \in \Domain {\OLId{latin-looped}{M}}$
شيء من المجال، لا رمز من اللغة؛ فلا تكون $!B({\OLId{latin-ordinary}{m}})$~!!a{formula}،
أي إن $\Subst{!B}{{\OLId{latin-ordinary}{m}}}{{\OLId{latin-ordinary}{x}}}$ غير معرّفة. وليس لنا أن نفترض وجود
!!{constant}s أو حدود تسمّي كل !!{element} من~$\Struct{M}$؛
فتعريف ال!!{structure}s لا يوجب ذلك. بل إن مجموعة
ال!!{constant}s في اللغة القياسية !!{denumerable}؛ فإن كان
$\Domain{{\OLId{latin-looped}{M}}}$ غير !!{enumerable} لم تكفِ ال!!{constant}s
لتسمية جميع الأشياء.

والخلاص من ذلك بإسنادات ال!!{variable}s، إذ نربط بها المتغيرات
مباشرة ب!!{element}s المجال. فبدلًا من أن نقول مثلًا إن
\iftag{prvEx}{$\lexists[{\OLId{latin-ordinary}{x}}][!B({\OLId{latin-ordinary}{x}})]$}{$\lforall[{\OLId{latin-ordinary}{x}}][!B({\OLId{latin-ordinary}{x}})]$} مُرضاة
في~$\Struct M$ إذا وفقط إذا
\iftag{prvEx}{كان $\Sat{{\OLId{latin-looped}{M}}}{!B({\OLId{latin-ordinary}{m}})}$ لواحد على الأقل من ${\OLId{latin-ordinary}{m}} \in
\Domain{{\OLId{latin-looped}{M}}}$}{كان ذلك لكل ${\OLId{latin-ordinary}{m}} \in \Domain{{\OLId{latin-looped}{M}}}$، $\Sat{{\OLId{latin-looped}{M}}}{!B({\OLId{latin-ordinary}{m}})}$}،
نقول: هي مُرضاة في~$\Struct M$ \emph{بالنسبة إلى}~${\OLId{latin-ordinary}{s}}$ إذا وفقط
إذا أُرضيت $!B({\OLId{latin-ordinary}{x}})$ بالنسبة إلى~$\Subst{{\OLId{latin-ordinary}{s}}}{{\OLId{latin-ordinary}{m}}}{{\OLId{latin-ordinary}{x}}}$
\iftag{prvEx}{لواحد على الأقل}{لكل} ${\OLId{latin-ordinary}{m}} \in \Domain {\OLId{latin-looped}{M}}$.
\end{explain}

\begin{ex}
لتكن $\Lang{L} = \{{\OLId{latin-ordinary}{a}}, {\OLId{latin-ordinary}{b}}, {\OLId{latin-ordinary}{f}}, {\OLId{latin-looped}{R}}\}$، حيث ${\OLId{latin-ordinary}{a}}$ و${\OLId{latin-ordinary}{b}}$ !!{constant}s،
و${\OLId{latin-ordinary}{f}}$~!!{function} ثنائية المحل، و${\OLId{latin-looped}{R}}$~!!{predicate} ثنائي المحل. وانظر
في ال!!{structure}~$\Struct{M}$ المعرّفة بما يأتي:
\begin{enumerate}
\item $\Domain {\OLId{latin-looped}{M}} = \{{\OLMathNumeral{1}}, {\OLMathNumeral{2}}, {\OLMathNumeral{3}}, {\OLMathNumeral{4}}\}$
\item $\Assign{{\OLId{latin-ordinary}{a}}}{{\OLId{latin-looped}{M}}} = {\OLMathNumeral{1}}$
\item $\Assign{{\OLId{latin-ordinary}{b}}}{{\OLId{latin-looped}{M}}} = {\OLMathNumeral{2}}$
\item $\Assign{{\OLId{latin-ordinary}{f}}}{{\OLId{latin-looped}{M}}}({\OLId{latin-ordinary}{x}}, {\OLId{latin-ordinary}{y}}) = {\OLId{latin-ordinary}{x}}+{\OLId{latin-ordinary}{y}}$ إذا كان ${\OLId{latin-ordinary}{x}}+{\OLId{latin-ordinary}{y}} \le {\OLMathNumeral{3}}$، و$= {\OLMathNumeral{3}}$ فيما عدا ذلك.
\item $\Assign{{\OLId{latin-looped}{R}}}{{\OLId{latin-looped}{M}}} = \{\tuple{{\OLMathNumeral{1}}, {\OLMathNumeral{1}}}, \tuple{{\OLMathNumeral{1}}, {\OLMathNumeral{2}}}, \tuple{{\OLMathNumeral{2}}, {\OLMathNumeral{3}}}, \tuple{{\OLMathNumeral{2}}, {\OLMathNumeral{4}}}\}$
\end{enumerate}
الدالة ${\OLId{latin-ordinary}{s}}({\OLId{latin-ordinary}{x}}) = {\OLMathNumeral{1}}$ التي تسند ${\OLMathNumeral{1}} \in \Domain{{\OLId{latin-looped}{M}}}$ إلى كل
!!{variable} إسناد متغيرات لـ~$\Struct{M}$.

عندئذ
\begin{align*}
\Value{{\OLId{latin-ordinary}{f}}({\OLId{latin-ordinary}{a}},{\OLId{latin-ordinary}{b}})}{{\OLId{latin-looped}{M}}}[{\OLId{latin-ordinary}{s}}] & = \Assign{{\OLId{latin-ordinary}{f}}}{{\OLId{latin-looped}{M}}}(\Value{{\OLId{latin-ordinary}{a}}}{{\OLId{latin-looped}{M}}}[{\OLId{latin-ordinary}{s}}], \Value{{\OLId{latin-ordinary}{b}}}{{\OLId{latin-looped}{M}}}[{\OLId{latin-ordinary}{s}}]).
\intertext{ولما كان ${\OLId{latin-ordinary}{a}}$ و${\OLId{latin-ordinary}{b}}$ !!{constant}s، فإن $\Value{{\OLId{latin-ordinary}{a}}}{{\OLId{latin-looped}{M}}}[{\OLId{latin-ordinary}{s}}]
  = \Assign{{\OLId{latin-ordinary}{a}}}{{\OLId{latin-looped}{M}}} = {\OLMathNumeral{1}}$ و$\Value{{\OLId{latin-ordinary}{b}}}{{\OLId{latin-looped}{M}}}[{\OLId{latin-ordinary}{s}}] = \Assign{{\OLId{latin-ordinary}{b}}}{{\OLId{latin-looped}{M}}} = {\OLMathNumeral{2}}$. ومن ثم}
\Value{{\OLId{latin-ordinary}{f}}({\OLId{latin-ordinary}{a}},{\OLId{latin-ordinary}{b}})}{{\OLId{latin-looped}{M}}}[{\OLId{latin-ordinary}{s}}] & = \Assign{{\OLId{latin-ordinary}{f}}}{{\OLId{latin-looped}{M}}}({\OLMathNumeral{1}}, {\OLMathNumeral{2}}) = {\OLMathNumeral{1}}+{\OLMathNumeral{2}} = {\OLMathNumeral{3}}.
\intertext{ولحساب قيمة ${\OLId{latin-ordinary}{f}}({\OLId{latin-ordinary}{f}}({\OLId{latin-ordinary}{a}},{\OLId{latin-ordinary}{b}}),{\OLId{latin-ordinary}{a}})$ يلزمنا النظر في}
\Value{{\OLId{latin-ordinary}{f}}({\OLId{latin-ordinary}{f}}({\OLId{latin-ordinary}{a}},{\OLId{latin-ordinary}{b}}),{\OLId{latin-ordinary}{a}})}{{\OLId{latin-looped}{M}}}[{\OLId{latin-ordinary}{s}}] & = \Assign{{\OLId{latin-ordinary}{f}}}{{\OLId{latin-looped}{M}}}(\Value{{\OLId{latin-ordinary}{f}}({\OLId{latin-ordinary}{a}}, {\OLId{latin-ordinary}{b}})}{{\OLId{latin-looped}{M}}}[{\OLId{latin-ordinary}{s}}],
\Value{{\OLId{latin-ordinary}{a}}}{{\OLId{latin-looped}{M}}}[{\OLId{latin-ordinary}{s}}]) = \Assign{{\OLId{latin-ordinary}{f}}}{{\OLId{latin-looped}{M}}}({\OLMathNumeral{3}}, {\OLMathNumeral{1}}) = {\OLMathNumeral{3}},
\intertext{لأن ${\OLMathNumeral{3}}+{\OLMathNumeral{1}} > {\OLMathNumeral{3}}$. ولما كان ${\OLId{latin-ordinary}{s}}({\OLId{latin-ordinary}{x}}) = {\OLMathNumeral{1}}$ و$\Value{{\OLId{latin-ordinary}{x}}}{{\OLId{latin-looped}{M}}}[{\OLId{latin-ordinary}{s}}] =
  {\OLId{latin-ordinary}{s}}({\OLId{latin-ordinary}{x}})$، كان لدينا أيضًا}
\Value{{\OLId{latin-ordinary}{f}}({\OLId{latin-ordinary}{f}}({\OLId{latin-ordinary}{a}},{\OLId{latin-ordinary}{b}}),{\OLId{latin-ordinary}{x}})}{{\OLId{latin-looped}{M}}}[{\OLId{latin-ordinary}{s}}] & = \Assign{{\OLId{latin-ordinary}{f}}}{{\OLId{latin-looped}{M}}}(\Value{{\OLId{latin-ordinary}{f}}({\OLId{latin-ordinary}{a}}, {\OLId{latin-ordinary}{b}})}{{\OLId{latin-looped}{M}}}[{\OLId{latin-ordinary}{s}}],
\Value{{\OLId{latin-ordinary}{x}}}{{\OLId{latin-looped}{M}}}[{\OLId{latin-ordinary}{s}}]) = \Assign{{\OLId{latin-ordinary}{f}}}{{\OLId{latin-looped}{M}}}({\OLMathNumeral{3}}, {\OLMathNumeral{1}}) = {\OLMathNumeral{3}},
\end{align*}

أما ال!!{formula} الذرية~${\OLId{latin-looped}{R}}({\OLId{latin-ordinary}{t}}_{\OLMathNumeral{1}}, {\OLId{latin-ordinary}{t}}_{\OLMathNumeral{2}})$، فتُرضى متى كانت
ثنائية قيمتي حجتيها، أي $\tuple{\Value{{\OLId{latin-ordinary}{t}}_{\OLMathNumeral{1}}}{{\OLId{latin-looped}{M}}}[{\OLId{latin-ordinary}{s}}], \Value{{\OLId{latin-ordinary}{t}}_{\OLMathNumeral{2}}}{{\OLId{latin-looped}{M}}}[{\OLId{latin-ordinary}{s}}]}$،
!!a{element} من~$\Assign{{\OLId{latin-looped}{R}}}{{\OLId{latin-looped}{M}}}$. فمن ذلك
$\Sat{{\OLId{latin-looped}{M}}}{{\OLId{latin-looped}{R}}({\OLId{latin-ordinary}{b}},{\OLId{latin-ordinary}{f}}({\OLId{latin-ordinary}{a}},{\OLId{latin-ordinary}{b}}))}[{\OLId{latin-ordinary}{s}}]$؛ إذ $\tuple{\Value{{\OLId{latin-ordinary}{b}}}{{\OLId{latin-looped}{M}}},
  \Value{{\OLId{latin-ordinary}{f}}({\OLId{latin-ordinary}{a}},{\OLId{latin-ordinary}{b}})}{{\OLId{latin-looped}{M}}}} = \tuple{{\OLMathNumeral{2}}, {\OLMathNumeral{3}}} \in \Assign{{\OLId{latin-looped}{R}}}{{\OLId{latin-looped}{M}}}$.
وأما $\Sat/{{\OLId{latin-looped}{M}}}{{\OLId{latin-looped}{R}}({\OLId{latin-ordinary}{x}}, {\OLId{latin-ordinary}{f}}({\OLId{latin-ordinary}{a}},{\OLId{latin-ordinary}{b}}))}[{\OLId{latin-ordinary}{s}}]$، فلأن
$\tuple{{\OLMathNumeral{1}}, {\OLMathNumeral{3}}} \notin \Assign{{\OLId{latin-looped}{R}}}{{\OLId{latin-looped}{M}}}$.

وأما الصيغة غير الذرية~$!A$، فسبيل الحكم عليها أن نأخذ من
التعريف الاستقرائي البند الموافق لرابطها الرئيس. فالرابط الرئيس
في~${\OLId{latin-looped}{R}}({\OLId{latin-ordinary}{a}}, {\OLId{latin-ordinary}{a}}) \lif ({\OLId{latin-looped}{R}}({\OLId{latin-ordinary}{b}}, {\OLId{latin-ordinary}{x}}) \lor {\OLId{latin-looped}{R}}({\OLId{latin-ordinary}{x}}, {\OLId{latin-ordinary}{b}}))$ هو~$\lif$؛ ولذلك
\begin{align*}
  & \Sat{{\OLId{latin-looped}{M}}}{{\OLId{latin-looped}{R}}({\OLId{latin-ordinary}{a}}, {\OLId{latin-ordinary}{a}}) \lif ({\OLId{latin-looped}{R}}({\OLId{latin-ordinary}{b}}, {\OLId{latin-ordinary}{x}}) \lor {\OLId{latin-looped}{R}}({\OLId{latin-ordinary}{x}}, {\OLId{latin-ordinary}{b}}))}[{\OLId{latin-ordinary}{s}}] \text{ إذا وفقط إذا }\\
  & \qquad
  \Sat/{{\OLId{latin-looped}{M}}}{{\OLId{latin-looped}{R}}({\OLId{latin-ordinary}{a}},{\OLId{latin-ordinary}{a}})}[{\OLId{latin-ordinary}{s}}] \text{ أو } \Sat{{\OLId{latin-looped}{M}}}{{\OLId{latin-looped}{R}}({\OLId{latin-ordinary}{b}}, {\OLId{latin-ordinary}{x}}) \lor {\OLId{latin-looped}{R}}({\OLId{latin-ordinary}{x}}, {\OLId{latin-ordinary}{b}})}[{\OLId{latin-ordinary}{s}}]
  \intertext{ولما كان $\Sat{{\OLId{latin-looped}{M}}}{{\OLId{latin-looped}{R}}({\OLId{latin-ordinary}{a}},{\OLId{latin-ordinary}{a}})}[{\OLId{latin-ordinary}{s}}]$ (لأن $\tuple{{\OLMathNumeral{1}},{\OLMathNumeral{1}}} \in
    \Assign{{\OLId{latin-looped}{R}}}{{\OLId{latin-looped}{M}}}$)، فلا يمكننا بعد تعيين الجواب، ويلزم أولًا أن
    نعرف هل $\Sat{{\OLId{latin-looped}{M}}}{{\OLId{latin-looped}{R}}({\OLId{latin-ordinary}{b}}, {\OLId{latin-ordinary}{x}}) \lor {\OLId{latin-looped}{R}}({\OLId{latin-ordinary}{x}}, {\OLId{latin-ordinary}{b}})}[{\OLId{latin-ordinary}{s}}]$:}
  & \Sat{{\OLId{latin-looped}{M}}}{{\OLId{latin-looped}{R}}({\OLId{latin-ordinary}{b}}, {\OLId{latin-ordinary}{x}}) \lor {\OLId{latin-looped}{R}}({\OLId{latin-ordinary}{x}}, {\OLId{latin-ordinary}{b}})}[{\OLId{latin-ordinary}{s}}] \text{ إذا وفقط إذا }\\
  & \qquad \Sat{{\OLId{latin-looped}{M}}}{{\OLId{latin-looped}{R}}({\OLId{latin-ordinary}{b}},{\OLId{latin-ordinary}{x}})}[{\OLId{latin-ordinary}{s}}] \text{ أو } \Sat{{\OLId{latin-looped}{M}}}{{\OLId{latin-looped}{R}}({\OLId{latin-ordinary}{x}}, {\OLId{latin-ordinary}{b}})}[{\OLId{latin-ordinary}{s}}]
  \intertext{وهذا هو الواقع، لأن $\Sat{{\OLId{latin-looped}{M}}}{{\OLId{latin-looped}{R}}({\OLId{latin-ordinary}{x}}, {\OLId{latin-ordinary}{b}})}[{\OLId{latin-ordinary}{s}}]$
    (إذ $\tuple{{\OLMathNumeral{1}},{\OLMathNumeral{2}}} \in \Assign{{\OLId{latin-looped}{R}}}{{\OLId{latin-looped}{M}}}$).}
\end{align*}

قد تقدم أن المغاير بالنسبة إلى~${\OLId{latin-ordinary}{x}}$ لـ~${\OLId{latin-ordinary}{s}}$ لا يجوز أن يخالف
الإسناد~${\OLId{latin-ordinary}{s}}$ إلا في قيمة~${\OLId{latin-ordinary}{x}}$، وقد يوافقه فيها أيضًا.
فلكل !!{element} من~$\Domain{{\OLId{latin-looped}{M}}}$ مغاير بالنسبة إلى~${\OLId{latin-ordinary}{x}}$ لـ~${\OLId{latin-ordinary}{s}}$:
\begin{align*}
  {\OLId{latin-ordinary}{s}}_{\OLMathNumeral{1}} & = \Subst{{\OLId{latin-ordinary}{s}}}{{\OLMathNumeral{1}}}{{\OLId{latin-ordinary}{x}}}, &
  {\OLId{latin-ordinary}{s}}_{\OLMathNumeral{2}} & = \Subst{{\OLId{latin-ordinary}{s}}}{{\OLMathNumeral{2}}}{{\OLId{latin-ordinary}{x}}},\\
  {\OLId{latin-ordinary}{s}}_{\OLMathNumeral{3}} & = \Subst{{\OLId{latin-ordinary}{s}}}{{\OLMathNumeral{3}}}{{\OLId{latin-ordinary}{x}}}, &
  {\OLId{latin-ordinary}{s}}_{\OLMathNumeral{4}} & = \Subst{{\OLId{latin-ordinary}{s}}}{{\OLMathNumeral{4}}}{{\OLId{latin-ordinary}{x}}}.
\end{align*}
فمثلًا، ${\OLId{latin-ordinary}{s}}_{\OLMathNumeral{2}}({\OLId{latin-ordinary}{x}}) = {\OLMathNumeral{2}}$ و${\OLId{latin-ordinary}{s}}_{\OLMathNumeral{2}}({\OLId{latin-ordinary}{y}}) = {\OLId{latin-ordinary}{s}}({\OLId{latin-ordinary}{y}}) = {\OLMathNumeral{1}}$ لجميع المتغيرات~${\OLId{latin-ordinary}{y}}$
غير~${\OLId{latin-ordinary}{x}}$. وهذه هي جميع المغايرات بالنسبة إلى~${\OLId{latin-ordinary}{x}}$ لـ~${\OLId{latin-ordinary}{s}}$ في
البنية~$\Struct{M}$، لأن $\Domain{{\OLId{latin-looped}{M}}} = \{{\OLMathNumeral{1}}, {\OLMathNumeral{2}}, {\OLMathNumeral{3}}, {\OLMathNumeral{4}}\}$. ولاحظ،
على وجه الخصوص، أن ${\OLId{latin-ordinary}{s}}_{\OLMathNumeral{1}} = {\OLId{latin-ordinary}{s}}$ (إذ يكون~${\OLId{latin-ordinary}{s}}$ دائمًا مغايرًا بالنسبة
إلى~${\OLId{latin-ordinary}{x}}$ لنفسه).

\iftag{prvEx}{لتعيين ما إذا كانت !!{formula} مكممة وجوديًا
~$\lexists[{\OLId{latin-ordinary}{x}}][!A({\OLId{latin-ordinary}{x}})]$ مُرضاة، يلزم أن نعيّن ما إذا كان
$\Sat{{\OLId{latin-looped}{M}}}{!A({\OLId{latin-ordinary}{x}})}[\Subst{{\OLId{latin-ordinary}{s}}}{{\OLId{latin-ordinary}{m}}}{{\OLId{latin-ordinary}{x}}}]$ لواحد على الأقل من
${\OLId{latin-ordinary}{m}} \in \Domain {\OLId{latin-looped}{M}}$. ومن ثم
  \[
  \Sat{{\OLId{latin-looped}{M}}}{\lexists[{\OLId{latin-ordinary}{x}}][({\OLId{latin-looped}{R}}({\OLId{latin-ordinary}{b}},{\OLId{latin-ordinary}{x}}) \lor {\OLId{latin-looped}{R}}({\OLId{latin-ordinary}{x}},{\OLId{latin-ordinary}{b}}))]}[{\OLId{latin-ordinary}{s}}],
  \]
لأن $\Sat{{\OLId{latin-looped}{M}}}{{\OLId{latin-looped}{R}}({\OLId{latin-ordinary}{b}},{\OLId{latin-ordinary}{x}}) \lor {\OLId{latin-looped}{R}}({\OLId{latin-ordinary}{x}}, {\OLId{latin-ordinary}{b}})}[\Subst{{\OLId{latin-ordinary}{s}}}{{\OLMathNumeral{1}}}{{\OLId{latin-ordinary}{x}}}]$
($\Subst{{\OLId{latin-ordinary}{s}}}{{\OLMathNumeral{3}}}{{\OLId{latin-ordinary}{x}}}$ يصلح أيضًا). لكن
  \[
  \Sat/{{\OLId{latin-looped}{M}}}{\lexists[{\OLId{latin-ordinary}{x}}][({\OLId{latin-looped}{R}}({\OLId{latin-ordinary}{b}},{\OLId{latin-ordinary}{x}}) \land {\OLId{latin-looped}{R}}({\OLId{latin-ordinary}{x}},{\OLId{latin-ordinary}{b}}))]}[{\OLId{latin-ordinary}{s}}]
  \]
لأنه أيًا كان ${\OLId{latin-ordinary}{m}} \in \Domain{{\OLId{latin-looped}{M}}}$ الذي نختاره، يكون
$\Sat/{{\OLId{latin-looped}{M}}}{{\OLId{latin-looped}{R}}({\OLId{latin-ordinary}{b}},{\OLId{latin-ordinary}{x}}) \land {\OLId{latin-looped}{R}}({\OLId{latin-ordinary}{x}},{\OLId{latin-ordinary}{b}})}[\Subst{{\OLId{latin-ordinary}{s}}}{{\OLId{latin-ordinary}{m}}}{{\OLId{latin-ordinary}{x}}}]$.}{}

\iftag{prvAll}{لتعيين ما إذا كانت !!{formula} مكممة كليًا
~$\lforall[{\OLId{latin-ordinary}{x}}][!A({\OLId{latin-ordinary}{x}})]$ مُرضاة، يلزم أن نعيّن ما إذا كان
$\Sat{{\OLId{latin-looped}{M}}}{!A({\OLId{latin-ordinary}{x}})}[\Subst{{\OLId{latin-ordinary}{s}}}{{\OLId{latin-ordinary}{m}}}{{\OLId{latin-ordinary}{x}}}]$ لكل ${\OLId{latin-ordinary}{m}} \in \Domain {\OLId{latin-looped}{M}}$. ومن ثم
  \[
  \Sat{{\OLId{latin-looped}{M}}}{\lforall[{\OLId{latin-ordinary}{x}}][({\OLId{latin-looped}{R}}({\OLId{latin-ordinary}{x}},{\OLId{latin-ordinary}{a}}) \lif {\OLId{latin-looped}{R}}({\OLId{latin-ordinary}{a}},{\OLId{latin-ordinary}{x}}))]}[{\OLId{latin-ordinary}{s}}],
  \]
لأن $\Sat{{\OLId{latin-looped}{M}}}{{\OLId{latin-looped}{R}}({\OLId{latin-ordinary}{x}},{\OLId{latin-ordinary}{a}}) \lif {\OLId{latin-looped}{R}}({\OLId{latin-ordinary}{a}},{\OLId{latin-ordinary}{x}})}[\Subst{{\OLId{latin-ordinary}{s}}}{{\OLId{latin-ordinary}{m}}}{{\OLId{latin-ordinary}{x}}}]$ لكل ${\OLId{latin-ordinary}{m}} \in
\Domain {\OLId{latin-looped}{M}}$. فعندما ${\OLId{latin-ordinary}{m}} = {\OLMathNumeral{1}}$ يكون
$\Sat{{\OLId{latin-looped}{M}}}{{\OLId{latin-looped}{R}}({\OLId{latin-ordinary}{a}},{\OLId{latin-ordinary}{x}})}[\Subst{{\OLId{latin-ordinary}{s}}}{{\OLMathNumeral{1}}}{{\OLId{latin-ordinary}{x}}}]$، فتكون التالية صادقة؛ وعندما
${\OLId{latin-ordinary}{m}} = {\OLMathNumeral{2}}$، أو ${\OLMathNumeral{3}}$، أو~${\OLMathNumeral{4}}$ يكون
$\Sat/{{\OLId{latin-looped}{M}}}{{\OLId{latin-looped}{R}}({\OLId{latin-ordinary}{x}},{\OLId{latin-ordinary}{a}})}[\Subst{{\OLId{latin-ordinary}{s}}}{{\OLId{latin-ordinary}{m}}}{{\OLId{latin-ordinary}{x}}}]$، فتكون المقدمة كاذبة. لكن
  \[
  \Sat/{{\OLId{latin-looped}{M}}}{\lforall[{\OLId{latin-ordinary}{x}}][({\OLId{latin-looped}{R}}({\OLId{latin-ordinary}{a}},{\OLId{latin-ordinary}{x}}) \lif {\OLId{latin-looped}{R}}({\OLId{latin-ordinary}{x}},{\OLId{latin-ordinary}{a}}))]}[{\OLId{latin-ordinary}{s}}]
  \]
لأن $\Sat/{{\OLId{latin-looped}{M}}}{{\OLId{latin-looped}{R}}({\OLId{latin-ordinary}{a}},{\OLId{latin-ordinary}{x}}) \lif {\OLId{latin-looped}{R}}({\OLId{latin-ordinary}{x}},{\OLId{latin-ordinary}{a}})}[\Subst{{\OLId{latin-ordinary}{s}}}{{\OLMathNumeral{2}}}{{\OLId{latin-ordinary}{x}}}]$ (إذ
$\Sat{{\OLId{latin-looped}{M}}}{{\OLId{latin-looped}{R}}({\OLId{latin-ordinary}{a}}, {\OLId{latin-ordinary}{x}})}[\Subst{{\OLId{latin-ordinary}{s}}}{{\OLMathNumeral{2}}}{{\OLId{latin-ordinary}{x}}}]$ و$\Sat/{{\OLId{latin-looped}{M}}}{{\OLId{latin-looped}{R}}({\OLId{latin-ordinary}{x}},
  {\OLId{latin-ordinary}{a}})}[\Subst{{\OLId{latin-ordinary}{s}}}{{\OLMathNumeral{2}}}{{\OLId{latin-ordinary}{x}}}]$).}{}

\iftag{defEx}{لتعيين ما إذا كانت !!{formula} مكممة وجوديًا
~$\lexists[{\OLId{latin-ordinary}{x}}][!A({\OLId{latin-ordinary}{x}})]$ مُرضاة، يلزم أن نعيّن ما إذا كان
$\Sat{{\OLId{latin-looped}{M}}}{\lnot \lforall[{\OLId{latin-ordinary}{x}}][\lnot !A({\OLId{latin-ordinary}{x}})]}[{\OLId{latin-ordinary}{s}}]$. فمثلًا، لدينا
  \[
  \Sat{{\OLId{latin-looped}{M}}}{\lexists[{\OLId{latin-ordinary}{x}}][({\OLId{latin-looped}{R}}({\OLId{latin-ordinary}{b}},{\OLId{latin-ordinary}{x}}) \lor {\OLId{latin-looped}{R}}({\OLId{latin-ordinary}{x}},{\OLId{latin-ordinary}{b}}))]}[{\OLId{latin-ordinary}{s}}].
  \]
أولًا، $\Sat{{\OLId{latin-looped}{M}}}{{\OLId{latin-looped}{R}}({\OLId{latin-ordinary}{b}},{\OLId{latin-ordinary}{x}}) \lor {\OLId{latin-looped}{R}}({\OLId{latin-ordinary}{x}}, {\OLId{latin-ordinary}{b}})}[\Subst{{\OLId{latin-ordinary}{s}}}{{\OLMathNumeral{1}}}{{\OLId{latin-ordinary}{x}}}]$
($\Subst{{\OLId{latin-ordinary}{s}}}{{\OLMathNumeral{3}}}{{\OLId{latin-ordinary}{x}}}$ يصلح أيضًا). ومن ثم
$\Sat/{{\OLId{latin-looped}{M}}}{\lnot({\OLId{latin-looped}{R}}({\OLId{latin-ordinary}{b}},{\OLId{latin-ordinary}{x}}) \lor {\OLId{latin-looped}{R}}({\OLId{latin-ordinary}{x}},{\OLId{latin-ordinary}{b}}))}[\Subst{{\OLId{latin-ordinary}{s}}}{{\OLMathNumeral{1}}}{{\OLId{latin-ordinary}{x}}}]$، وبالتالي
$\Sat/{{\OLId{latin-looped}{M}}}{\lforall[{\OLId{latin-ordinary}{x}}][\lnot(({\OLId{latin-looped}{R}}({\OLId{latin-ordinary}{b}},{\OLId{latin-ordinary}{x}}) \lor {\OLId{latin-looped}{R}}({\OLId{latin-ordinary}{x}},{\OLId{latin-ordinary}{b}})))]}[{\OLId{latin-ordinary}{s}}]$، ومنه
$\Sat{{\OLId{latin-looped}{M}}}{\lnot\lforall[{\OLId{latin-ordinary}{x}}][\lnot(({\OLId{latin-looped}{R}}({\OLId{latin-ordinary}{b}},{\OLId{latin-ordinary}{x}}) \lor
{\OLId{latin-looped}{R}}({\OLId{latin-ordinary}{x}},{\OLId{latin-ordinary}{b}})))]}[{\OLId{latin-ordinary}{s}}]$. وأما من الجهة الأخرى
  \[
  \Sat/{{\OLId{latin-looped}{M}}}{\lexists[{\OLId{latin-ordinary}{x}}][({\OLId{latin-looped}{R}}({\OLId{latin-ordinary}{b}},{\OLId{latin-ordinary}{x}}) \land {\OLId{latin-looped}{R}}({\OLId{latin-ordinary}{x}},{\OLId{latin-ordinary}{b}}))]}[{\OLId{latin-ordinary}{s}}].
  \]
وذلك لأن $\Sat{{\OLId{latin-looped}{M}}}{\lforall[{\OLId{latin-ordinary}{x}}][\lnot({\OLId{latin-looped}{R}}({\OLId{latin-ordinary}{b}},{\OLId{latin-ordinary}{x}}) \land
{\OLId{latin-looped}{R}}({\OLId{latin-ordinary}{x}},{\OLId{latin-ordinary}{b}}))]}[{\OLId{latin-ordinary}{s}}]$، إذ لا يوجد ${\OLId{latin-ordinary}{m}} \in \Domain {\OLId{latin-looped}{M}}$ بحيث
$\Sat{{\OLId{latin-looped}{M}}}{{\OLId{latin-looped}{R}}({\OLId{latin-ordinary}{b}},{\OLId{latin-ordinary}{x}}) \land {\OLId{latin-looped}{R}}({\OLId{latin-ordinary}{x}},{\OLId{latin-ordinary}{b}})}[\Subst{{\OLId{latin-ordinary}{s}}}{{\OLId{latin-ordinary}{m}}}{{\OLId{latin-ordinary}{x}}}]$. ويُستدل من هذه الأمثلة على القاعدة الآتية: يكون $\Sat{{\OLId{latin-looped}{M}}}{\lexists[{\OLId{latin-ordinary}{x}}][!A({\OLId{latin-ordinary}{x}})]}[{\OLId{latin-ordinary}{s}}]$ إذا
وفقط إذا كان $\Sat{{\OLId{latin-looped}{M}}}{!A({\OLId{latin-ordinary}{x}})}[\Subst{{\OLId{latin-ordinary}{s}}}{{\OLId{latin-ordinary}{m}}}{{\OLId{latin-ordinary}{x}}}]$ لواحد على الأقل
من~${\OLId{latin-ordinary}{m}} \in \Domain {\OLId{latin-looped}{M}}$.}{}

\iftag{defAll}{لتعيين ما إذا كانت !!{formula} مكممة كليًا
~$\lforall[{\OLId{latin-ordinary}{x}}][!A({\OLId{latin-ordinary}{x}})]$ مُرضاة، يلزم أن نعيّن ما إذا كان
$\Sat{{\OLId{latin-looped}{M}}}{\lnot\lexists[{\OLId{latin-ordinary}{x}}][\lnot !A({\OLId{latin-ordinary}{x}})]}[{\OLId{latin-ordinary}{s}}]$. فمثلًا
  \[
  \Sat{{\OLId{latin-looped}{M}}}{\lforall[{\OLId{latin-ordinary}{x}}][({\OLId{latin-looped}{R}}({\OLId{latin-ordinary}{x}},{\OLId{latin-ordinary}{a}}) \lif {\OLId{latin-looped}{R}}({\OLId{latin-ordinary}{a}},{\OLId{latin-ordinary}{x}}))]}[{\OLId{latin-ordinary}{s}}],
  \]
أولًا، $\Sat{{\OLId{latin-looped}{M}}}{{\OLId{latin-looped}{R}}({\OLId{latin-ordinary}{x}},{\OLId{latin-ordinary}{a}}) \lif {\OLId{latin-looped}{R}}({\OLId{latin-ordinary}{a}},{\OLId{latin-ordinary}{x}})}[\Subst{{\OLId{latin-ordinary}{s}}}{{\OLId{latin-ordinary}{m}}}{{\OLId{latin-ordinary}{x}}}]$ لكل
${\OLId{latin-ordinary}{m}} \in \Domain {\OLId{latin-looped}{M}}$ ($\Sat{{\OLId{latin-looped}{M}}}{{\OLId{latin-looped}{R}}({\OLId{latin-ordinary}{a}},{\OLId{latin-ordinary}{x}})}[\Subst{{\OLId{latin-ordinary}{s}}}{{\OLMathNumeral{1}}}{{\OLId{latin-ordinary}{x}}}]$
و$\Sat/{{\OLId{latin-looped}{M}}}{{\OLId{latin-looped}{R}}({\OLId{latin-ordinary}{a}},{\OLId{latin-ordinary}{x}})}[\Subst{{\OLId{latin-ordinary}{s}}}{{\OLId{latin-ordinary}{m}}}{{\OLId{latin-ordinary}{x}}}]$ عندما ${\OLId{latin-ordinary}{m}} = {\OLMathNumeral{2}}$، أو ${\OLMathNumeral{3}}$،
أو~${\OLMathNumeral{4}}$). ومن ثم لا يوجد ${\OLId{latin-ordinary}{m}} \in \Domain {\OLId{latin-looped}{M}}$ بحيث
$\Sat{{\OLId{latin-looped}{M}}}{\lnot({\OLId{latin-looped}{R}}({\OLId{latin-ordinary}{x}},{\OLId{latin-ordinary}{a}}) \lif {\OLId{latin-looped}{R}}({\OLId{latin-ordinary}{a}},{\OLId{latin-ordinary}{x}}))}[\Subst{{\OLId{latin-ordinary}{s}}}{{\OLId{latin-ordinary}{m}}}{{\OLId{latin-ordinary}{x}}}]$، ولذلك
$\Sat/{{\OLId{latin-looped}{M}}}{\lexists[{\OLId{latin-ordinary}{x}}][\lnot({\OLId{latin-looped}{R}}({\OLId{latin-ordinary}{x}},{\OLId{latin-ordinary}{a}}) \lif {\OLId{latin-looped}{R}}({\OLId{latin-ordinary}{a}},{\OLId{latin-ordinary}{x}}))]}[{\OLId{latin-ordinary}{s}}]$. وبالتالي
$\Sat{{\OLId{latin-looped}{M}}}{\lnot\lexists[{\OLId{latin-ordinary}{x}}][\lnot({\OLId{latin-looped}{R}}({\OLId{latin-ordinary}{x}},{\OLId{latin-ordinary}{a}}) \lif {\OLId{latin-looped}{R}}({\OLId{latin-ordinary}{a}},{\OLId{latin-ordinary}{x}}))]}[{\OLId{latin-ordinary}{s}}]$. ومن
جهة أخرى
  \[
  \Sat/{{\OLId{latin-looped}{M}}}{\lforall[{\OLId{latin-ordinary}{x}}][({\OLId{latin-looped}{R}}({\OLId{latin-ordinary}{a}},{\OLId{latin-ordinary}{x}}) \lif {\OLId{latin-looped}{R}}({\OLId{latin-ordinary}{x}},{\OLId{latin-ordinary}{a}}))]}[{\OLId{latin-ordinary}{s}}],
  \]
لأن $\Sat/{{\OLId{latin-looped}{M}}}{{\OLId{latin-looped}{R}}({\OLId{latin-ordinary}{a}},{\OLId{latin-ordinary}{x}}) \lif {\OLId{latin-looped}{R}}({\OLId{latin-ordinary}{x}},{\OLId{latin-ordinary}{a}})}[\Subst{{\OLId{latin-ordinary}{s}}}{{\OLMathNumeral{2}}}{{\OLId{latin-ordinary}{x}}}]$، ومن ثم
$\Sat{{\OLId{latin-looped}{M}}}{\lexists[{\OLId{latin-ordinary}{x}}][\lnot({\OLId{latin-looped}{R}}({\OLId{latin-ordinary}{a}},{\OLId{latin-ordinary}{x}}) \lif {\OLId{latin-looped}{R}}({\OLId{latin-ordinary}{x}},{\OLId{latin-ordinary}{a}}))]}[{\OLId{latin-ordinary}{s}}]$. وهذه الأمثلة تدل على القاعدة الآتية: يكون
$\Sat{{\OLId{latin-looped}{M}}}{\lforall[{\OLId{latin-ordinary}{x}}][!A({\OLId{latin-ordinary}{x}})]}[{\OLId{latin-ordinary}{s}}]$ إذا وفقط إذا كان
$\Sat{{\OLId{latin-looped}{M}}}{!A({\OLId{latin-ordinary}{x}})}[\Subst{{\OLId{latin-ordinary}{s}}}{{\OLId{latin-ordinary}{m}}}{{\OLId{latin-ordinary}{x}}}]$ لكل ${\OLId{latin-ordinary}{m}} \in \Domain {\OLId{latin-looped}{M}}$.}{}

ولننظر الآن في الحالة الأكثر تركيبًا:
\[
\lforall[{\OLId{latin-ordinary}{x}}][({\OLId{latin-looped}{R}}({\OLId{latin-ordinary}{a}},{\OLId{latin-ordinary}{x}}) \lif \lexists[{\OLId{latin-ordinary}{y}}][{\OLId{latin-looped}{R}}({\OLId{latin-ordinary}{x}},{\OLId{latin-ordinary}{y}})])].
\]
لما كان $\Sat/{{\OLId{latin-looped}{M}}}{{\OLId{latin-looped}{R}}({\OLId{latin-ordinary}{a}},{\OLId{latin-ordinary}{x}})}[\Subst{{\OLId{latin-ordinary}{s}}}{{\OLMathNumeral{3}}}{{\OLId{latin-ordinary}{x}}}]$ و
$\Sat/{{\OLId{latin-looped}{M}}}{{\OLId{latin-looped}{R}}({\OLId{latin-ordinary}{a}},{\OLId{latin-ordinary}{x}})}[\Subst{{\OLId{latin-ordinary}{s}}}{{\OLMathNumeral{4}}}{{\OLId{latin-ordinary}{x}}}]$، فلم يبق للنظر في تالية الشرطية إلا الحالتان ${\OLId{latin-ordinary}{m}} = {\OLMathNumeral{1}}$
و${\OLId{latin-ordinary}{m}} = {\OLMathNumeral{2}}$. هل $\Sat{{\OLId{latin-looped}{M}}}{\lexists[{\OLId{latin-ordinary}{y}}][{\OLId{latin-looped}{R}}({\OLId{latin-ordinary}{x}},{\OLId{latin-ordinary}{y}})]}[\Subst{{\OLId{latin-ordinary}{s}}}{{\OLMathNumeral{1}}}{{\OLId{latin-ordinary}{x}}}]$؟
يكون الجواب بالإيجاب إذا وُجد ${\OLId{latin-ordinary}{n}} \in \Domain {\OLId{latin-looped}{M}}$ واحد على الأقل بحيث
$\Sat{{\OLId{latin-looped}{M}}}{{\OLId{latin-looped}{R}}({\OLId{latin-ordinary}{x}},{\OLId{latin-ordinary}{y}})}[\Subst{\Subst{{\OLId{latin-ordinary}{s}}}{{\OLMathNumeral{1}}}{{\OLId{latin-ordinary}{x}}}}{{\OLId{latin-ordinary}{n}}}{{\OLId{latin-ordinary}{y}}}]$. ولنا أن
نأخذ ${\OLId{latin-ordinary}{n}} = {\OLMathNumeral{1}}$، فيكون
$\Subst{\Subst{{\OLId{latin-ordinary}{s}}}{{\OLMathNumeral{1}}}{{\OLId{latin-ordinary}{x}}}}{{\OLId{latin-ordinary}{n}}}{{\OLId{latin-ordinary}{y}}} = \Subst{{\OLId{latin-ordinary}{s}}}{{\OLMathNumeral{1}}}{{\OLId{latin-ordinary}{y}}} = {\OLId{latin-ordinary}{s}}$. ولما كان
${\OLId{latin-ordinary}{s}}({\OLId{latin-ordinary}{x}}) = {\OLMathNumeral{1}}$، و${\OLId{latin-ordinary}{s}}({\OLId{latin-ordinary}{y}}) = {\OLMathNumeral{1}}$، و$\tuple{{\OLMathNumeral{1}},{\OLMathNumeral{1}}} \in \Assign{{\OLId{latin-looped}{R}}}{{\OLId{latin-looped}{M}}}$، فالجواب
نعم.

وأما السؤال هل
$\Sat{{\OLId{latin-looped}{M}}}{\lexists[{\OLId{latin-ordinary}{y}}][{\OLId{latin-looped}{R}}({\OLId{latin-ordinary}{x}},{\OLId{latin-ordinary}{y}})]}[\Subst{{\OLId{latin-ordinary}{s}}}{{\OLMathNumeral{2}}}{{\OLId{latin-ordinary}{x}}}]$، فجوابه بالنظر في
إسنادات ال!!{variable}s~$\Subst{\Subst{{\OLId{latin-ordinary}{s}}}{{\OLMathNumeral{2}}}{{\OLId{latin-ordinary}{x}}}}{{\OLId{latin-ordinary}{n}}}{{\OLId{latin-ordinary}{y}}}$. وهنا،
عندما ${\OLId{latin-ordinary}{n}} = {\OLMathNumeral{1}}$، يكون هذا الإسناد~${\OLId{latin-ordinary}{s}}_{\OLMathNumeral{2}} = \Subst{{\OLId{latin-ordinary}{s}}}{{\OLMathNumeral{2}}}{{\OLId{latin-ordinary}{x}}}$، وهذا الإسناد لا
يرضي ${\OLId{latin-looped}{R}}({\OLId{latin-ordinary}{x}},{\OLId{latin-ordinary}{y}})$ (${\OLId{latin-ordinary}{s}}_{\OLMathNumeral{2}}({\OLId{latin-ordinary}{x}}) = {\OLMathNumeral{2}}$، و${\OLId{latin-ordinary}{s}}_{\OLMathNumeral{2}}({\OLId{latin-ordinary}{y}}) = {\OLMathNumeral{1}}$،
و$\tuple{{\OLMathNumeral{2}},{\OLMathNumeral{1}}}\notin \Assign{{\OLId{latin-looped}{R}}}{{\OLId{latin-looped}{M}}}$). وأما إذا أخذنا
$\Subst{\Subst{{\OLId{latin-ordinary}{s}}}{{\OLMathNumeral{2}}}{{\OLId{latin-ordinary}{x}}}}{{\OLMathNumeral{3}}}{{\OLId{latin-ordinary}{y}}} = \Subst{{\OLId{latin-ordinary}{s}}_{\OLMathNumeral{2}}}{{\OLMathNumeral{3}}}{{\OLId{latin-ordinary}{y}}}$،
فلدينا $\Sat{{\OLId{latin-looped}{M}}}{{\OLId{latin-looped}{R}}({\OLId{latin-ordinary}{x}},{\OLId{latin-ordinary}{y}})}[\Subst{{\OLId{latin-ordinary}{s}}_{\OLMathNumeral{2}}}{{\OLMathNumeral{3}}}{{\OLId{latin-ordinary}{y}}}]$ لأن
$\tuple{{\OLMathNumeral{2}},{\OLMathNumeral{3}}} \in \Assign{{\OLId{latin-looped}{R}}}{{\OLId{latin-looped}{M}}}$، ومن ثم
$\Sat{{\OLId{latin-looped}{M}}}{\lexists[{\OLId{latin-ordinary}{y}}][{\OLId{latin-looped}{R}}({\OLId{latin-ordinary}{x}},{\OLId{latin-ordinary}{y}})]}[{\OLId{latin-ordinary}{s}}_{\OLMathNumeral{2}}]$.

وهكذا، لكل ${\OLId{latin-ordinary}{m}} \in \Domain {\OLId{latin-looped}{M}}$، إما
$\Sat/{{\OLId{latin-looped}{M}}}{{\OLId{latin-looped}{R}}({\OLId{latin-ordinary}{a}},{\OLId{latin-ordinary}{x}})}[\Subst{{\OLId{latin-ordinary}{s}}}{{\OLId{latin-ordinary}{m}}}{{\OLId{latin-ordinary}{x}}}]$ (إذا كان ${\OLId{latin-ordinary}{m}} = {\OLMathNumeral{3}}$، أو~${\OLMathNumeral{4}}$)،
وإما $\Sat{{\OLId{latin-looped}{M}}}{\lexists[{\OLId{latin-ordinary}{y}}][{\OLId{latin-looped}{R}}({\OLId{latin-ordinary}{x}},{\OLId{latin-ordinary}{y}})]}[\Subst{{\OLId{latin-ordinary}{s}}}{{\OLId{latin-ordinary}{m}}}{{\OLId{latin-ordinary}{x}}}]$ (إذا كان
${\OLId{latin-ordinary}{m}} = {\OLMathNumeral{1}}$، أو~${\OLMathNumeral{2}}$)، ومن ثم
\[
\Sat{{\OLId{latin-looped}{M}}}{\lforall[{\OLId{latin-ordinary}{x}}][({\OLId{latin-looped}{R}}({\OLId{latin-ordinary}{a}},{\OLId{latin-ordinary}{x}}) \lif \lexists[{\OLId{latin-ordinary}{y}}][{\OLId{latin-looped}{R}}({\OLId{latin-ordinary}{x}},{\OLId{latin-ordinary}{y}})])]}[{\OLId{latin-ordinary}{s}}].
\]
ومن جهة أخرى
\[
\Sat/{{\OLId{latin-looped}{M}}}{\lexists[{\OLId{latin-ordinary}{x}}][({\OLId{latin-looped}{R}}({\OLId{latin-ordinary}{a}},{\OLId{latin-ordinary}{x}}) \land \lforall[{\OLId{latin-ordinary}{y}}][{\OLId{latin-looped}{R}}({\OLId{latin-ordinary}{x}},{\OLId{latin-ordinary}{y}})])]}[{\OLId{latin-ordinary}{s}}].
\]
لدينا $\Sat{{\OLId{latin-looped}{M}}}{{\OLId{latin-looped}{R}}({\OLId{latin-ordinary}{a}},{\OLId{latin-ordinary}{x}})}[\Subst{{\OLId{latin-ordinary}{s}}}{{\OLId{latin-ordinary}{m}}}{{\OLId{latin-ordinary}{x}}}]$ فقط عندما ${\OLId{latin-ordinary}{m}} = {\OLMathNumeral{1}}$ وعندما
${\OLId{latin-ordinary}{m}} = {\OLMathNumeral{2}}$. لكن لكل من هاتين القيمتين لـ~${\OLId{latin-ordinary}{m}}$ يوجد بدوره
${\OLId{latin-ordinary}{n}} \in \Domain {\OLId{latin-looped}{M}}$، هو ${\OLId{latin-ordinary}{n}} = {\OLMathNumeral{4}}$، بحيث
$\Sat/{{\OLId{latin-looped}{M}}}{{\OLId{latin-looped}{R}}({\OLId{latin-ordinary}{x}},{\OLId{latin-ordinary}{y}})}[\Subst{\Subst{{\OLId{latin-ordinary}{s}}}{{\OLId{latin-ordinary}{m}}}{{\OLId{latin-ordinary}{x}}}}{{\OLId{latin-ordinary}{n}}}{{\OLId{latin-ordinary}{y}}}]$، ومن ثم
$\Sat/{{\OLId{latin-looped}{M}}}{\lforall[{\OLId{latin-ordinary}{y}}][{\OLId{latin-looped}{R}}({\OLId{latin-ordinary}{x}},{\OLId{latin-ordinary}{y}})]}[\Subst{{\OLId{latin-ordinary}{s}}}{{\OLId{latin-ordinary}{m}}}{{\OLId{latin-ordinary}{x}}}]$ عندما ${\OLId{latin-ordinary}{m}} = {\OLMathNumeral{1}}$
وعندما ${\OLId{latin-ordinary}{m}} = {\OLMathNumeral{2}}$. فالحاصل أنه لا يوجد ${\OLId{latin-ordinary}{m}} \in \Domain {\OLId{latin-looped}{M}}$ بحيث
$\Sat{{\OLId{latin-looped}{M}}}{{\OLId{latin-looped}{R}}({\OLId{latin-ordinary}{a}},{\OLId{latin-ordinary}{x}}) \land \lforall[{\OLId{latin-ordinary}{y}}][{\OLId{latin-looped}{R}}({\OLId{latin-ordinary}{x}},{\OLId{latin-ordinary}{y}})]}[\Subst{{\OLId{latin-ordinary}{s}}}{{\OLId{latin-ordinary}{m}}}{{\OLId{latin-ordinary}{x}}}]$.
\end{ex}

\iftag{defEx}{%
\begin{prop}\ollabel{prop:sat-ex}
$\Sat{{\OLId{latin-looped}{M}}}{\lexists[{\OLId{latin-ordinary}{x}}][!B({\OLId{latin-ordinary}{x}})]}[{\OLId{latin-ordinary}{s}}]$ إذا وفقط إذا وُجد مغاير بالنسبة
إلى~${\OLId{latin-ordinary}{x}}$، ${\OLId{latin-ordinary}{s}}'$، لـ~${\OLId{latin-ordinary}{s}}$ بحيث $\Sat{{\OLId{latin-looped}{M}}}{!B({\OLId{latin-ordinary}{x}})}[{\OLId{latin-ordinary}{s}}']$.
\end{prop}

\begin{proof}
  تمرين.
\end{proof}
}{}

\tagprob{defEx}
\begin{prob}
  أثبت \olref[fol][syn][sat]{prop:sat-ex}.
\end{prob}
\tagendprob

\iftag{defAll}{%
\begin{prop}\ollabel{prop:sat-all}
$\Sat{{\OLId{latin-looped}{M}}}{\lforall[{\OLId{latin-ordinary}{x}}][!B({\OLId{latin-ordinary}{x}})]}[{\OLId{latin-ordinary}{s}}]$ إذا وفقط إذا كان، لكل مغاير
بالنسبة إلى~${\OLId{latin-ordinary}{x}}$، ${\OLId{latin-ordinary}{s}}'$، لـ~${\OLId{latin-ordinary}{s}}$، $\Sat{{\OLId{latin-looped}{M}}}{!B({\OLId{latin-ordinary}{x}})}[{\OLId{latin-ordinary}{s}}']$.
\end{prop}

\begin{proof}
  تمرين.
\end{proof}
}{}

\tagprob{defAll}
\begin{prob}
  أثبت \olref[fol][syn][sat]{prop:sat-all}.
\end{prob}
\tagendprob

\begin{prob}
لتكن $\Lang L = \{{\OLId{latin-ordinary}{c}}, {\OLId{latin-ordinary}{f}}, {\OLId{latin-looped}{A}}\}$ ذات !!{constant} واحد، و!!{function}
واحدة أحادية المحل، و!!{predicate} واحد ثنائي المحل، ولتكن
ال!!{structure}~$\Struct{M}$ معطاة بما يأتي:
\begin{enumerate}
\item $\Domain {\OLId{latin-looped}{M}} = \{{\OLMathNumeral{1}}, {\OLMathNumeral{2}}, {\OLMathNumeral{3}}\}$
\item $\Assign{{\OLId{latin-ordinary}{c}}}{{\OLId{latin-looped}{M}}} = {\OLMathNumeral{3}}$
\item $\Assign{{\OLId{latin-ordinary}{f}}}{{\OLId{latin-looped}{M}}}({\OLMathNumeral{1}}) = {\OLMathNumeral{2}}, \Assign{{\OLId{latin-ordinary}{f}}}{{\OLId{latin-looped}{M}}}({\OLMathNumeral{2}}) = {\OLMathNumeral{3}}, \Assign{{\OLId{latin-ordinary}{f}}}{{\OLId{latin-looped}{M}}}({\OLMathNumeral{3}}) = {\OLMathNumeral{2}}$
\item $\Assign{{\OLId{latin-looped}{A}}}{{\OLId{latin-looped}{M}}} = \{\tuple{{\OLMathNumeral{1}}, {\OLMathNumeral{2}}}, \tuple{{\OLMathNumeral{2}}, {\OLMathNumeral{3}}}, \tuple{{\OLMathNumeral{3}}, {\OLMathNumeral{3}}}\}$
\end{enumerate}
(أ) ليكن ${\OLId{latin-ordinary}{s}}({\OLId{latin-ordinary}{v}}) = {\OLMathNumeral{1}}$ لكل ال!!{variable}s~${\OLId{latin-ordinary}{v}}$. عيّن ما إذا كان
\[
\Sat{{\OLId{latin-looped}{M}}}{\lexists[{\OLId{latin-ordinary}{x}}][({\OLId{latin-looped}{A}}({\OLId{latin-ordinary}{f}}({\OLId{latin-ordinary}{z}}), {\OLId{latin-ordinary}{c}}) \lif \lforall[{\OLId{latin-ordinary}{y}}][({\OLId{latin-looped}{A}}({\OLId{latin-ordinary}{y}}, {\OLId{latin-ordinary}{x}}) \lor {\OLId{latin-looped}{A}}({\OLId{latin-ordinary}{f}}({\OLId{latin-ordinary}{y}}),
      {\OLId{latin-ordinary}{x}}))])]}[{\OLId{latin-ordinary}{s}}]
\]
وبيّن وجه الحكم.

(ب) ضع بنية أخرى وإسنادًا آخر لل!!{variable}s بحيث لا تكون
ال!!{formula} مُرضاة بهما.
\end{prob}

\end{document}
